\import{test/algebra-spike/01-ring-reducts.tex}
\import{algebra/group/power.tex}
\import{algebra/ordered.tex}

\section{Monoid-power compatibility}

This fixture gives natural powers one monoid-level iteration owner and checks
that the existing group and ring terms retain their old equations under the
new bodies.

\begin{definition}\label{spike16_monoid_right_map}
    $\spikeMonoidRightMap{M}{x}=
        \{(y,\mul[M](y,x))\mid y\in\carrier[M]\}$.
\end{definition}

\begin{proposition}\label{spike16_monoid_right_map_iff}
    $(y,z)\in\spikeMonoidRightMap{M}{x}$ iff
        $y\in\carrier[M]$ and $z=\mul[M](y,x)$.
\end{proposition}
\begin{proof}
    Follows by
        \auto{spike16_monoid_right_map,pair_eq_iff}.
\end{proof}

\begin{proposition}\label{spike16_monoid_right_map_type}
    Let $M$ be a monoid.
    Let $x\in\carrier[M]$.
    Then $\spikeMonoidRightMap{M}{x}\in
        \funs{\carrier[M]}{\carrier[M]}$.
\end{proposition}
\begin{proof}
    Let $f=\spikeMonoidRightMap{M}{x}$.
    $f$ is a relation by
        \auto{relation,spike16_monoid_right_map}.
    $f$ is right-unique by
        \auto{rightunique,spike16_monoid_right_map_iff}.
    Hence $f$ is a function by \auto.
    $\dom{f}=\carrier[M]$ by
        \auto{setext,dom_iff,spike16_monoid_right_map_iff}.
    Show that for every $y\in\dom{f}$ we have $f(y)\in\carrier[M]$.
    \begin{subproof}
        Fix $y\in\dom{f}$.
        Hence $y\in\carrier[M]$ by \auto.
        $f(y)=\mul[M](y,x)$ by
            \auto{function_apply_intro,spike16_monoid_right_map_iff}.
        Follows by \auto{monoid_is_magma,magma_welldef}.
    \end{subproof}
    Hence $f$ is a function to $\carrier[M]$ by \auto.
    Follows by \auto{funs_intro}.
\end{proof}

\begin{proposition}\label{spike16_monoid_right_map_apply}
    Let $M$ be a monoid.
    Let $x,y\in\carrier[M]$.
    Then $\apply{\spikeMonoidRightMap{M}{x}}{y}=\mul[M](y,x)$.
\end{proposition}
\begin{proof}
    $\spikeMonoidRightMap{M}{x}$ is a function by
        \auto{spike16_monoid_right_map_type,funs_is_function}.
    $(y,\mul[M](y,x))\in\spikeMonoidRightMap{M}{x}$ by
        \auto{spike16_monoid_right_map_iff}.
    Follows by \auto{function_apply_intro}.
\end{proof}

\begin{definition}\label{spike16_monoid_power}
    $\spikeMonoidPower{M}{x}{n}=
        \iterationAt{\spikeMonoidRightMap{M}{x}}{\neutral[M]}{n}$.
\end{definition}

\begin{proposition}\label{spike16_monoid_power_zero}
    $\spikeMonoidPower{M}{x}{\emptyset}=\neutral[M]$.
\end{proposition}
\begin{proof}
    Follows by
        \auto{spike16_monoid_power,iteration_at_zero}.
\end{proof}

\begin{proposition}\label{spike16_monoid_power_type}
    Let $M$ be a monoid.
    Let $x\in\carrier[M]$.
    Let $n$ be a natural number.
    Then $\spikeMonoidPower{M}{x}{n}\in\carrier[M]$.
\end{proposition}
\begin{proof}
    $\spikeMonoidRightMap{M}{x}\in
        \funs{\carrier[M]}{\carrier[M]}$
        by \auto{spike16_monoid_right_map_type}.
    $\neutral[M]\in\carrier[M]$ by
        \auto{monoid_is_unital_magma,unitalmagma_type}.
    Follows by
        \auto{spike16_monoid_power,iteration_at_closed}.
\end{proof}

\begin{proposition}\label{spike16_monoid_power_successor}
    Let $M$ be a monoid.
    Let $x\in\carrier[M]$.
    Let $n$ be a natural number.
    Then $\spikeMonoidPower{M}{x}{\suc{n}}
        =\mul[M](\spikeMonoidPower{M}{x}{n},x)$.
\end{proposition}
\begin{proof}
    $\spikeMonoidPower{M}{x}{n}\in\carrier[M]$
        by \auto{spike16_monoid_power_type}.
    $\spikeMonoidPower{M}{x}{\suc{n}}
        =\apply{\spikeMonoidRightMap{M}{x}}{\spikeMonoidPower{M}{x}{n}}$
        by \auto{spike16_monoid_power,iteration_at_successor}.
    Follows by \auto{spike16_monoid_right_map_apply}.
\end{proof}

\subsection{The group compatibility view}

\begin{definition}\label{spike16_group_power_body}
    $\spikeGroupPower{G}{a}{n}=\spikeMonoidPower{G}{a}{n}$.
\end{definition}

\begin{proposition}\label{spike16_group_right_map_is_old_translation}
    $\spikeMonoidRightMap{G}{a}=\RightTranslation{G}{a}$.
\end{proposition}
\begin{proof}
    Follows by
        \auto{spike16_monoid_right_map,right_translation,setext}.
\end{proof}

\begin{theorem}\label{spike16_group_power_iteration_eq}
    $\spikeGroupPower{G}{a}{n}
        =\iterationAt{\RightTranslation{G}{a}}{\neutral[G]}{n}$.
\end{theorem}
\begin{proof}
    Follows by
        \auto{spike16_group_power_body,spike16_monoid_power,
            spike16_group_right_map_is_old_translation}.
\end{proof}

\begin{theorem}\label{spike16_group_power_new_body_eq_old}
    $\spikeGroupPower{G}{a}{n}=\groupPower{G}{a}{n}$.
\end{theorem}
\begin{proof}
    Follows by
        \auto{spike16_group_power_iteration_eq,group_power}.
\end{proof}

\subsection{The ring-natural compatibility view}

\begin{definition}\label{spike16_ring_natural_body}
    $\spikeRingNatural{R}{n}=
        \spikeMonoidPower{\AdditiveGroup{R}}{\ringOne[R]}{n}$.
\end{definition}

\begin{proposition}\label{spike16_ring_natural_zero}
    $\spikeRingNatural{R}{\emptyset}=\ringZero[R]$.
\end{proposition}
\begin{proof}
    Follows by
        \auto{spike16_ring_natural_body,spike16_monoid_power_zero,
            spike_additive_group_fields}.
\end{proof}

\begin{proposition}\label{spike16_ring_natural_successor}
    Let $R$ be a ring.
    Let $n$ be a natural number.
    Then $\spikeRingNatural{R}{\suc{n}}
        =\ringAdd[R](\spikeRingNatural{R}{n},\ringOne[R])$.
\end{proposition}
\begin{proof}
    $\AdditiveGroup{R}$ is a monoid by
        \auto{spike_additive_group_is_exact,commutative_group,group}.
    $\ringOne[R]\in\carrier[R]$ by \auto{ring_one_type}.
    Hence $\ringOne[R]\in\carrier[\AdditiveGroup{R}]$
        by \auto{spike_additive_group_fields}.
    $\spikeRingNatural{R}{\suc{n}}
        =\mul[\AdditiveGroup{R}](
            \spikeRingNatural{R}{n},\ringOne[R])$
        by \auto{spike16_ring_natural_body,
            spike16_monoid_power_successor}.
    Follows by \auto{spike_additive_group_fields}.
\end{proof}

\begin{theorem}\label{spike16_ring_natural_new_body_eq_old}
    Let $R$ be a ring.
    Let $n$ be a natural number.
    Then $\spikeRingNatural{R}{n}=\ringNatural{R}{n}$.
\end{theorem}
\begin{proof}
    Let $C=\{k\in\naturals\mid
        \spikeRingNatural{R}{k}=\ringNatural{R}{k}\}$.
    $\emptyset\in C$ by
        \auto{omega_empty,spike16_ring_natural_zero,
            ring_natural_zero}.
    Show that for every $k\in C$ we have $\suc{k}\in C$.
    \begin{subproof}
        Fix $k\in C$.
        Then $k$ is a natural number and
            $\spikeRingNatural{R}{k}=\ringNatural{R}{k}$.
        $\spikeRingNatural{R}{\suc{k}}
            =\ringAdd[R](\spikeRingNatural{R}{k},\ringOne[R])$
            by \auto{spike16_ring_natural_successor}.
        $\ringNatural{R}{\suc{k}}
            =\ringAdd[R](\ringNatural{R}{k},\ringOne[R])$
            by \auto{ring_natural_successor}.
        $\suc{k}$ is a natural number by \auto{omega_successor}.
        Follows by \auto.
    \end{subproof}
    Hence $C$ is an inductive set by \auto{inductive_set}.
    Hence every natural number is an element of $C$
        by \auto{omega_minimal}.
    Thus $n\in C$.
    Follows by \auto.
\end{proof}

\begin{theorem}\label{spike16_ring_natural_recursion_eq}
    Let $R$ be a ring.
    Let $n$ be a natural number.
    Then $\spikeRingNatural{R}{n}
        =\Apply{\naturalRecOp}{\ringZero[R],
            \Apply{\ringNaturalStepOp}{R},n}$.
\end{theorem}
\begin{proof}
    Follows by
        \auto{spike16_ring_natural_new_body_eq_old,ring_natural}.
\end{proof}
